\documentclass[12pt, twoside]{article}
\input{/Users/chris/github/course-files/docs/ib/preamble}
\setlength{\parskip}{0.3\baselineskip}

\fancyhead[LE]{\thepage}
\fancyhead[RO]{Markscheme}
\fancyhead[LO]{La Scuola d'Italia / Dr.\ Huson / IB Math 12 HL \\* 14 October 2026}

\begin{document}
\subsubsection*{12.6 Quiz: Complex Numbers --- Markscheme}

\emph{Calculator allowed. Where a part asks for an exact form, a GDC
decimal earns the method marks shown but not the final A1.}

\begin{enumerate}[itemsep=0.3cm]

%--- Problem 1: Cartesian form; conjugate; product and quotient; modulus and argument [6 marks] ---
\item[1.] [Maximum mark: 6]

$z_{1} = 3 + 4i$, \quad $z_{2} = 1 - 2i$.

\textbf{(a)} $z_{1}z_{2}^{*} = (3 + 4i)(1 + 2i) = 3 + 6i + 4i + 8i^{2}$
\hfill \textbf{M1}

$= -5 + 10i$ \hfill \textbf{A1}

\emph{M1 for using $z_{2}^{*} = 1 + 2i$ and expanding. Using $z_{2}$ in
place of its conjugate, giving $11 - 2i$, earns M0A0.}

\textbf{(b)} $\dfrac{3 + 4i}{1 - 2i} \cdot \dfrac{1 + 2i}{1 + 2i}
= \dfrac{-5 + 10i}{5}$ \hfill \textbf{M1}

$= -1 + 2i$ \hfill \textbf{A1}

\emph{M1 for multiplying numerator and denominator by $1 + 2i$ (the
numerator is the product from (a)), or for a GDC entry shown. FT: award M1
for dividing the candidate's answer to (a) by $|z_{2}|^{2} = 5$.}

\textbf{(c)} $\left|\dfrac{z_{1}}{z_{2}}\right| = \sqrt{(-1)^{2} + 2^{2}}
= \sqrt{5} \; (= 2.24)$ \hfill \textbf{A1}

$\arg\dfrac{z_{1}}{z_{2}} = \pi - \arctan 2 = 2.0344\ldots = 2.03$
\hfill \textbf{A1}

\emph{Accept $\frac{|z_{1}|}{|z_{2}|} = \frac{5}{\sqrt{5}}$ for the
modulus. For the argument accept $2.03$ radians or $117^{\circ}$
($116.565\ldots^{\circ}$). $\arctan(-2) = -1.11$ (the fourth-quadrant
angle) earns A0: $-1 + 2i$ lies in the second quadrant. FT from the
candidate's quotient in (b).}

\newpage

%--- Problem 2: Euler and polar form; product in polar and Cartesian form; geometric meaning of multiplication [7 marks] ---
\item[2.] [Maximum mark: 7]

$z = -\sqrt{3} + i$, \quad
$w = 3\left(\cos\frac{\pi}{4} + i\sin\frac{\pi}{4}\right)$.

\textbf{(a)} $|z| = \sqrt{3 + 1} = 2$ \hfill \textbf{A1}

$\arg z = \pi - \dfrac{\pi}{6} = \dfrac{5\pi}{6}$, so
$z = 2e^{i\frac{5\pi}{6}}$ \hfill \textbf{A1}

\emph{Second quadrant: $\arctan\left(-\frac{1}{\sqrt{3}}\right) =
-\frac{\pi}{6}$ alone earns A0 for the argument. Accept $2.62$ for the
argument.}

\textbf{(b)} $zw = 2 \times 3 \operatorname{cis}\left(\dfrac{5\pi}{6} +
\dfrac{\pi}{4}\right) = 6\operatorname{cis}\dfrac{13\pi}{12}$
\hfill \textbf{M1}

$= 6\left(\cos\left(-\dfrac{11\pi}{12}\right) +
i\sin\left(-\dfrac{11\pi}{12}\right)\right)$ \hfill \textbf{A1}

\emph{M1 for multiplying the moduli and adding the arguments. A1 requires
the argument in the stated range; $\frac{13\pi}{12}$ left unadjusted earns
M1A0. Accept $-2.88$ for the argument. FT from (a).}

\textbf{(c)} $zw = -5.80 - 1.55i$ \hfill \textbf{A1}

\emph{Exact form: $zw = -\frac{3}{2}(\sqrt{6} + \sqrt{2}) -
\frac{3}{2}(\sqrt{6} - \sqrt{2})\,i = -5.7955\ldots - 1.5529\ldots i$.
Accept either the exact or the 3 s.f.\ form, from (b) or from multiplying
the Cartesian forms directly.}

\textbf{(d)} A rotation about the origin through $\dfrac{\pi}{4}$
anticlockwise \hfill \textbf{A1}

and an enlargement, centre the origin, scale factor $3$ \hfill \textbf{A1}

\emph{The two marks are independent. The rotation needs the angle and the
direction; ``turned by $45^{\circ}$ anticlockwise'' is fine. ``Stretched by
$3$'' or ``three times as far from the origin'' earns the second A1.}

\newpage

%--- Problem 3: Conjugate root theorem; sum and product of roots of a cubic; Argand triangle [13 marks] ---
\item[3.] [Maximum mark: 13]

$z^{3} + az^{2} + bz + 26 = 0$, \quad $a, b \in \mathbb{R}$, root
$3 + 2i$.

\textbf{(a)} $3 - 2i$ \hfill \textbf{A1}

The coefficients are real, so non-real roots occur in conjugate pairs.
\hfill \textbf{R1}

\emph{R1 requires the reference to real coefficients (or to $a$, $b$ and
$26$ being real).}

\textbf{(b)} Product of the roots $= -26$ \hfill \textbf{M1}

$(3 + 2i)(3 - 2i)\,\gamma = 13\gamma$ \hfill \textbf{A1}

$13\gamma = -26$, so $\gamma = -2$ \hfill \textbf{A1}

\emph{M1 for $(-1)^{3} \times 26$ or for $\alpha\beta\gamma = -26$, stated
or used. The product $+26$ leads to $\gamma = 2$: M0A1A0. FT from the
candidate's second root in (a).}

\textbf{(c)} \emph{Method 1 (sum and pairwise products):}

Sum of the roots $= (3 + 2i) + (3 - 2i) + (-2) = 4 = -a$
\hfill \textbf{M1}

$a = -4$ \hfill \textbf{A1}

Sum of pairwise products $= 13 + (-2)(3 + 2i) + (-2)(3 - 2i)
= 13 - 12 = 1$ \hfill \textbf{M1}

$b = 1$ \hfill \textbf{A1}

\emph{Method 2 (expanding the factors):}

$(z + 2)(z^{2} - 6z + 13)$ \hfill \textbf{M1}

$= z^{3} - 4z^{2} + z + 26$ \hfill \textbf{M1}

$a = -4$, \quad $b = 1$ \hfill \textbf{A1A1}

\emph{Substituting $z = 3 + 2i$ and equating real and imaginary parts is
also valid (M1 for the substitution, M1 for equating both parts). FT from
the candidate's roots.}

\textbf{(d)} A $(3, 2)$, \; B $(3, -2)$, \; C $(-2, 0)$

$\text{AC} = |(3 + 2i) - (-2)| = |5 + 2i| = \sqrt{29}$ and
$\text{BC} = |5 - 2i| = \sqrt{29}$ \hfill \textbf{M1}

$\text{AC} = \text{BC}$, so triangle ABC is isosceles. \hfill \textbf{A1 AG}

\emph{M1 for both lengths found (as moduli or by Pythagoras). Accept the
symmetry argument: A and B are reflections of each other in the real axis,
on which C lies, so $\text{CA} = \text{CB}$.}

\textbf{(e)} Base $\text{AB} = 4$, height $= 3 - (-2) = 5$
\hfill \textbf{M1}

$\text{Area} = \frac{1}{2} \times 4 \times 5 = 10$ \hfill \textbf{A1}

\emph{FT from the candidate's real root.}

\newpage

%--- Problem 4: Proof of de Moivre by induction; multiple-angle identity; exact power [14 marks] ---
\item[4.] [Maximum mark: 14]

\textbf{(a)} Let $P(n)$ be
$(\cos\theta + i\sin\theta)^{n} = \cos n\theta + i\sin n\theta$.

$n = 1$: $(\cos\theta + i\sin\theta)^{1} = \cos\theta + i\sin\theta$, so
$P(1)$ is true. \hfill \textbf{A1}

Assume $P(k)$ is true for some $k \in \mathbb{Z}^{+}$:
$(\cos\theta + i\sin\theta)^{k} = \cos k\theta + i\sin k\theta$.
\hfill \textbf{M1}

$(\cos\theta + i\sin\theta)^{k+1}
= (\cos k\theta + i\sin k\theta)(\cos\theta + i\sin\theta)$
\hfill \textbf{M1}

$= \cos k\theta\cos\theta - \sin k\theta\sin\theta
+ i(\sin k\theta\cos\theta + \cos k\theta\sin\theta)$
\hfill \textbf{A1}

$= \cos(k + 1)\theta + i\sin(k + 1)\theta$ \hfill \textbf{A1}

$P(1)$ is true, and $P(k)$ true implies $P(k + 1)$ true; so $P(n)$ is true
for all $n \in \mathbb{Z}^{+}$. \hfill \textbf{R1}

\emph{The assumption must be stated as an assumption (``assume'',
``suppose'', ``let''); ``$P(k)$ is true'' stated as fact earns M0. The
second A1 needs the compound angle identities applied to both parts. R1
only if the first four marks are earned and the conclusion names both the
base case and the inductive step.}

\textbf{(b)} $(\cos\theta + i\sin\theta)^{4} = \cos 4\theta + i\sin 4\theta$
\hfill \textbf{M1}

Write $c = \cos\theta$, $s = \sin\theta$:

$(c + is)^{4} = c^{4} + 4ic^{3}s - 6c^{2}s^{2} - 4ics^{3} + s^{4}$
\hfill \textbf{M1}

Real parts: $\cos 4\theta = c^{4} - 6c^{2}s^{2} + s^{4}$ \hfill \textbf{A1}

$= c^{4} - 6c^{2}(1 - c^{2}) + (1 - c^{2})^{2}$ \hfill \textbf{M1}

$= c^{4} - 6c^{2} + 6c^{4} + 1 - 2c^{2} + c^{4}
= 8c^{4} - 8c^{2} + 1$ \hfill \textbf{A1 AG}

\emph{Second M1 for a binomial expansion with all five terms (powers of $i$
need not yet be simplified). A1 for the real part with signs correct.
Third M1 for replacing $s^{2}$ by $1 - c^{2}$ throughout. The final A1 needs
the expansion of $(1 - c^{2})^{2}$ seen.}

\textbf{(c)} $1 - \sqrt{3}\,i = 2e^{-i\frac{\pi}{3}}$ \hfill \textbf{A1}

$\left(1 - \sqrt{3}\,i\right)^{10} = 2^{10}e^{-i\frac{10\pi}{3}}
= 1024e^{i\frac{2\pi}{3}}$ \hfill \textbf{M1}

$= 1024\left(-\dfrac{1}{2} + \dfrac{\sqrt{3}}{2}i\right)
= -512 + 512\sqrt{3}\,i$ \hfill \textbf{A1}

\emph{A1 for both modulus $2$ and argument $-\frac{\pi}{3}$ (accept
$\frac{5\pi}{3}$). M1 for raising the modulus to the $10$th power and
multiplying the argument by $10$. The GDC value $-512 + 886.81i$ alone
earns A0M0A0: the question asks for de Moivre's theorem and the form
$a + b\sqrt{3}\,i$; it is a useful check.}

\newpage

%--- Problem 5: Euler form; cube roots (rational exponent); sum of roots; w^(2/3) [13 marks] ---
\item[5.] [Maximum mark: 13]

$w = -4\sqrt{2} + 4\sqrt{2}\,i$.

\textbf{(a)} $|w| = \sqrt{32 + 32} = 8$ \hfill \textbf{A1}

$\arg w = \dfrac{3\pi}{4}$, so $w = 8e^{i\frac{3\pi}{4}}$ \hfill \textbf{A1}

\emph{Accept $2.36$ for the argument.}

\textbf{(b)} $z = 8^{\frac{1}{3}}e^{i\frac{1}{3}\left(\frac{3\pi}{4}
+ 2k\pi\right)}$ \hfill \textbf{M1}

$z_{1} = 2e^{i\frac{\pi}{4}}$ \hfill \textbf{A1}

Successive roots differ in argument by $\dfrac{2\pi}{3}$
\hfill \textbf{M1}

$z_{2} = 2e^{i\frac{11\pi}{12}}$ \hfill \textbf{A1}

$z_{3} = 2e^{-i\frac{5\pi}{12}}$ \hfill \textbf{A1}

\emph{First M1 for the cube root of the modulus with the argument divided
by $3$ (de Moivre with $n = \frac{1}{3}$). $z_{3}$ given as
$2e^{i\frac{19\pi}{12}}$ (outside the stated range) earns A0 for that root
only. FT from (a).}

\textbf{(c)} $2\left(\cos\dfrac{\pi}{4} + i\sin\dfrac{\pi}{4}\right)
= \sqrt{2} + \sqrt{2}\,i$ \hfill \textbf{A1}

\textbf{(d)} \emph{Method 1:} $z^{3} - w = 0$ has no $z^{2}$ term
\hfill \textbf{M1}

so the sum of the roots $= -\dfrac{0}{1} = 0$ \hfill \textbf{A1 AG}

\emph{Method 2:} $z_{1} + z_{2} + z_{3}
= z_{1}\left(1 + e^{i\frac{2\pi}{3}} + e^{i\frac{4\pi}{3}}\right)$
\hfill \textbf{M1}

$= z_{1}\left(1 + \left(-\dfrac{1}{2} + \dfrac{\sqrt{3}}{2}i\right)
+ \left(-\dfrac{1}{2} - \dfrac{\sqrt{3}}{2}i\right)\right) = 0$
\hfill \textbf{A1 AG}

\emph{Accept adding the three roots in Cartesian form (exactly). A GDC sum
alone, with no structure shown, earns M1A0.}

\textbf{(e)} $w^{\frac{2}{3}} = \left(w^{\frac{1}{3}}\right)^{2}$, so
square each cube root \hfill \textbf{M1}

$\left(2e^{i\frac{\pi}{4}}\right)^{2} = 4e^{i\frac{\pi}{2}} = 4i$
\hfill \textbf{A1}

$\left(2e^{i\frac{11\pi}{12}}\right)^{2} = 4e^{-i\frac{\pi}{6}}
= 2\sqrt{3} - 2i$, \quad
$\left(2e^{-i\frac{5\pi}{12}}\right)^{2} = 4e^{-i\frac{5\pi}{6}}
= -2\sqrt{3} - 2i$ \hfill \textbf{A1}

\emph{Accept $\left(w^{2}\right)^{\frac{1}{3}}$ with $w^{2} = 64e^{i\frac{3\pi}{2}}$.
Accept $\pm 3.46 - 2i$. The final A1 needs both of the last two values. A
single value ($4i$) with no others earns M1A1A0. FT from (b).}

\end{enumerate}
\end{document}
